\documentclass[11pt]{article}
\usepackage[T1]{fontenc}
\usepackage[margin=1in]{geometry}
\usepackage{amsmath,amssymb,amsthm}
\usepackage{hyperref}
\setlength{\emergencystretch}{3em}
\newtheorem{theorem}{Theorem}
\newtheorem{lemma}[theorem]{Lemma}
\newtheorem{proposition}[theorem]{Proposition}
\newcommand{\Z}{\mathbb Z}
\newcommand{\N}{\mathbb N}
\newcommand{\R}{\mathbb R}
\newcommand{\abs}[1]{\lvert#1\rvert}
\newcommand{\pos}[1]{(#1)_{+}}
\title{Square-uniform Fibonacci minima:\\a complete formal proof for one interpretation of MathDB 371202}
\author{Myeongjin Kang}
\date{26 October 2026}
\begin{document}
\maketitle
\begin{abstract}
For the bi-infinite integer Fibonacci recurrence, initialized uniformly on the integer square $[-N,N]^2$, we prove that the probability of a global absolute minimum at zero tends to $5/24$. The result holds both for a unique minimum and for a minimum with ties allowed. The proof includes an exact characterization of the minimum event, an exact lattice count, and the limiting argument. All these bridges are checked in Lean 4.34.1. The mathematical $5/24$ argument is due to Shivam Patel's MathDB reply of 20 August 2026. The primary source's Conjecture 20 states $P(0)=1/4$ but does not specify square-uniform sampling. Thus our result is a counterexample to the expanding-square interpretation, not a distribution-independent resolution of that underspecified conjecture.
\end{abstract}

\noindent\textit{Credit and license.} This AI-assisted paper verifies Patel's existing argument. Original prose is CC BY 4.0 and original code is MIT.

\section{Source, scope, and attribution}
Marc T. Pudelko's \emph{Modular Periodicity of Random Initialized Recurrences}, arXiv:2510.24882v2 (3 November 2025), Conjecture 20, equation (20), proposes a distribution for the location of the absolute minimum of a Fibonacci recurrence and its chiral recurrence, including $P(0)=1/4$ \cite{pudelko}. The MathDB entry numbered 371202 records that conjecture \cite{mathdb}.

There is an essential qualification. In the primary source, the phrase describing the initialization specifies integer initial values, but gives neither a finite sampling set nor an exhaustion procedure or probability law on $\Z^2$. In particular, the source text does \emph{not} explicitly impose the expanding-square convention used below. There is no uniform probability measure on all of $\Z^2$ assigning the same mass to each point. A probability law, or a limiting convention, must therefore be supplied before $P(0)$ has a precise meaning.

Patel's reply on MathDB, dated 20 August 2026, supplies the expanding-square counterexample, its cone criterion, its lattice sum, and the value $5/24$ \cite{patel}. We credit that mathematical argument to Patel. His reply also says that square-uniform initialization is explicitly specified by the source; our direct check of arXiv v2 does not support that assertion. We preserve the argument while correcting its source attribution. We make no new-to-literature claim, no claim that the unspecified probability in the source must mean square-uniform probability, and no claim about the complete probability distribution at other indices.

A failure already for the Fibonacci recurrence suffices to invalidate a formula asserted for both recurrences under a convention that includes this recurrence. We do not need, or claim here, a separate formal analysis of the chiral recurrence. We also do not count the separate MathDB 371201 weight-preservation formalization as part of this outcome.

\section{Precise model and main result}
For $x,y\in\Z$, let $a^{x,y}\colon\Z\to\Z$ be the solution with
\[
 a_0=x,\qquad a_1=y,\qquad a_{n+2}=a_n+a_{n+1}\quad(n\in\Z).
\]
Both directions are determined: the inverse step is $a_n=a_{n+2}-a_{n+1}$. For $N\in\N$, put
\[
 S_N=[-N,N]^2\cap\Z^2,\qquad \#S_N=(2N+1)^2.
\]
All points of $S_N$ have equal probability. Define
\begin{align*}
 U_N&=\{(x,y)\in S_N:\abs{a_0}<\abs{a_n}\text{ for every }n\ne0\},\\
 W_N&=\{(x,y)\in S_N:\abs{a_0}\le\abs{a_n}\text{ for every }n\in\Z\}.
\end{align*}
Thus $U_N$ is the unique-minimum event, while $W_N$ permits ties. Let $p_N=\#U_N/(2N+1)^2$ and $w_N=\#W_N/(2N+1)^2$. We use $\pos{u}=\max(u,0)$ for an integer $u$.

\begin{theorem}\label{thm:main}
For $k\ge1$ set
\[
 T_k(N)=\sum_{j=1}^{N}\pos{N-kj},\qquad q_k=\left\lfloor\frac{N}{k}\right\rfloor.
\]
Then
\begin{align}
 \#U_N&=2N+2T_2(N)+2T_3(N),\label{eq:count}\\
 2T_k(N)&=q_k\bigl(2N-k(q_k+1)\bigr),\label{eq:triangle}\\
 \#U_N&\le\#W_N\le\#U_N+2(2N+1).\label{eq:tie}
\end{align}
Consequently,
\[
 \lim_{N\to\infty}p_N=\lim_{N\to\infty}w_N=\frac5{24}\ne\frac14.
\]
In particular the square-uniform interpretation of $P(0)=1/4$ is false, whether or not tied minima are included.
\end{theorem}

\section{The complete minimum criterion}
The five neighboring terms are
\[
 a_{-2}=2x-y,\quad a_{-1}=y-x,\quad a_0=x,\quad a_1=y,\quad a_2=x+y.
\]
The following assertion includes the infinite tails; it is not merely a check of these five terms.

\begin{lemma}\label{lem:cone}
Zero is the unique global absolute minimum if and only if one of the following disjoint alternatives holds:
\begin{enumerate}
\item $x=0$ and $y\ne0$;
\item $x>0$ and either $y<-2x$ or $y>3x$;
\item $x<0$ and either $y>-2x$ or $y<3x$.
\end{enumerate}
Equivalently, for $x\ne0$, the criterion is $y/x<-2$ or $y/x>3$.
\end{lemma}
\begin{proof}
Simultaneous negation preserves all absolute values, so first suppose $x>0$. Necessity follows from the four strict inequalities
\[
 x<\abs y,\quad x<\abs{y-x},\quad
 x<\abs{x+y},\quad x<\abs{2x-y}.
\]
For $y<0$, the first inequality gives $y<-x$; then $x+y<0$, and the third inequality gives $y<-2x$. For $y\ge0$, the first gives $y>x$; the second then gives $y>2x$. Hence $2x-y<0$, and the fourth gives $y>3x$.

For sufficiency, the elementary tail facts are these. Two positive terms each larger than a nonnegative threshold $t$ generate, by addition, further positive terms larger than $t$. The same holds for two negative terms after simultaneously negating them. For the backward recursion
\[
 b_{n+2}=b_n-b_{n+1},
\]
two consecutive terms of opposite signs with absolute values greater than $t$ continue with alternating signs and absolute values greater than $t$: subtracting a term of the opposite sign adds their absolute values. Each fact follows by induction, and applies to every index in the corresponding infinite tail.

If $y>3x$, the forward terms $a_1=y$ and $a_2=x+y$ are positive and exceed $x$. The backward terms $a_{-1}=y-x$ and $a_{-2}=2x-y$ have opposite signs and absolute values greater than $x$. Thus all terms away from zero have absolute value greater than $x$.

If $y<-2x$, the forward terms $y$ and $x+y$ are negative and have absolute values greater than $x$. The backward terms $y-x$ and $2x-y$ again have opposite signs and absolute values greater than $x$. The same inductions apply.

When $x=0$ and $y\ne0$, use the same argument with threshold zero, choosing the sign according to $y$. When $(x,y)=(0,0)$ the recurrence is identically zero and its minimum is not unique. For $x<0$, simultaneous negation reduces the claim to $-x>0$ and reverses the displayed inequalities as stated.
\end{proof}

\section{Exact square-lattice counting}
\begin{proof}[Proof of \eqref{eq:count} and \eqref{eq:triangle}]
The column $x=0$ contributes exactly $2N$ points, namely all nonzero $y$ in $[-N,N]$. Fix $x=j>0$. By Lemma~\ref{lem:cone}, the permissible integer $y$ form the two disjoint intervals
\[
 [-N,-2j-1]\cap\Z,\qquad [3j+1,N]\cap\Z.
\]
Their cardinalities are $\pos{N-2j}$ and $\pos{N-3j}$ respectively. The negative columns have the same cardinalities by $(x,y)\mapsto(-x,-y)$. Summing all columns proves \eqref{eq:count}. In particular, this count is the cardinality of the actual global-minimum event, not an independently defined proxy count.

For $k\ge1$, terms with $j>q_k$ are zero, and terms with $1\le j\le q_k$ are nonnegative. Therefore
\[
 T_k(N)=\sum_{j=1}^{q_k}(N-kj)
       =q_kN-k\frac{q_k(q_k+1)}2.
\]
Multiplying by two gives \eqref{eq:triangle}. If $k\mid N$, including the final term $j=N/k$ simply includes a zero. This accounts for the harmless difference between this formula and Patel's sum truncated at $\lfloor(N-1)/k\rfloor$. The formulas are valid also at $N=0$.
\end{proof}

\section{Limits and the treatment of ties}
\begin{proof}[Proof of the limit for unique minima]
Write $N=kq_k+r_k$ with $0\le r_k<k$. For fixed $k\ge1$ and $N>0$,
\[
 \frac{q_k}{N}=\frac1k\left(1-\frac{r_k}{N}\right)
 \longrightarrow\frac1k,
\]
since $0\le r_k/N\le k/N\to0$. Dividing \eqref{eq:triangle} by $N^2$ gives
\[
 \frac{T_k(N)}{N^2}=
 \frac12\frac{q_k}{N}\left(2-k\left(\frac{q_k}{N}+\frac1N\right)\right)
 \longrightarrow\frac1{2k}.
\]
It follows from \eqref{eq:count} that
\[
 \frac{\#U_N}{N^2}\longrightarrow
 2\cdot\frac14+2\cdot\frac16=\frac56,
 \qquad
 \frac{(2N+1)^2}{N^2}=\left(2+\frac1N\right)^2\longrightarrow4.
\]
Taking the quotient yields $p_N\to(5/6)/4=5/24$.
\end{proof}

\begin{proof}[Proof of \eqref{eq:tie} and the limit with ties]
A weak minimum at zero satisfies the same four neighboring inequalities with $\le$ in place of $<$. For $x\ge0$, these imply
\[
 y\le-2x\quad\text{or}\quad y\ge3x.
\]
One can see this by replacing each assertion $x\le\abs z$ by the disjunction $x\le z$ or $x\le-z$ and intersecting the four conditions. Simultaneous negation gives the corresponding assertion when $x\le0$. Unless $y=-2x$ or $y=3x$, these alternatives are precisely the strict cones of Lemma~\ref{lem:cone}; when $x=0$ the only excluded case is $(0,0)$, which is already on both lines. Thus every point in $W_N\setminus U_N$ lies on one of the two lines
\[
 y=-2x,\qquad y=3x.
\]
Each line has at most $2N+1$ points in the square: its point is determined by its integer first coordinate. This proves \eqref{eq:tie}, without needing a convention assigning tied minima to one index.

For $N>0$, divide that bound by $(2N+1)^2$ to obtain
\[
 p_N\le w_N\le p_N+\frac{2}{2N+1}\le p_N+\frac2N.
\]
The squeeze theorem gives $w_N\to5/24$. Every rule which counts all unique minima at zero and only counts actual minima at zero gives probabilities between these two endpoints and hence has the same limit. Finally, $5/24\ne1/4$ by rational arithmetic, so uniqueness of limits excludes convergence to $1/4$.
\end{proof}

\section{Lean coverage and validation}
The single Lean module is \texttt{formal/FibonacciMinima.lean}, in namespace \texttt{FibonacciMinima}. Its definitions construct the positive and negative tails by two-step recursion, join them into an integer-indexed sequence, and define the finite square and its filtered minimum events directly. The following are the main proof interfaces:
\begin{center}
\begin{tabular}{ll}
Mathematical statement & Lean declaration\\\hline
Bi-infinite recurrence & \texttt{sequence\_recurrence}\\
Strict global-event bridge & \texttt{global\_minimum\_iff}\\
Exact cone criterion & \texttt{minimum\_iff\_cone}\\
Actual lattice cardinality & \texttt{exact\_count}\\
Exact truncated-sum evaluation & \texttt{triangle\_formula}\\
Quotient asymptotics & \texttt{quotient\_ratio\_limit}\\
Unique-minimum probability limit & \texttt{square\_uniform\_limit}\\
Tie boundary and count bound & \texttt{weak\_strict\_or\_boundary}, \texttt{weak\_count\_bounds}\\
Weak global-event bridge & \texttt{weak\_global\_minimum\_iff}\\
Tie-inclusive probability limit & \texttt{weak\_square\_uniform\_limit}
\end{tabular}
\end{center}
The bundled main theorem is \texttt{fibonacci\_minima\_square\_counterexample}. It has no hypotheses supplying an unproved recurrence, count, or limit bridge. It combines the recurrence, the global strict-minimum characterization, exact count, both limiting probabilities, and exclusion of $1/4$ for the tie-inclusive event.

The pinned verifier compiled the module and a separate declaration-audit module under Lean 4.34.1 with warnings treated as errors. The audit covers 48 declarations. Its observed foundational dependencies are exactly \texttt{propext}, \texttt{Classical.choice}, and \texttt{Quot.sound}. No unproved proof placeholders or additional postulates are present. Dependency revisions and the source hash are recorded in \texttt{build/audit.json} and the source-provenance record.

The analytic steps use ordinary real limits in mathlib: bounded remainders are squeezed to zero, then multiplication, addition, and division of limits are applied. This formal result establishes the entire expanding-square model above. It does not establish an interpretation of the original source's omitted probability law. It does not prove the Lucas-ratio law for other indices, classify all possible sampling laws, or claim that the arctangent expression is false for every law. The accompanying independent review assesses both the formalization and this explicitly restricted scope.

\begin{thebibliography}{9}
\bibitem{pudelko} Marc T. Pudelko, \emph{Modular Periodicity of Random Initialized Recurrences}, arXiv:2510.24882v2, 3 November 2025, Conjecture 20, equation (20), Remarks 21--22. \url{https://arxiv.org/html/2510.24882v2}.
\bibitem{mathdb} MathDB 371202, \emph{Probability distribution of Fibonacci sequence minima}, \url{https://mathdb.com/p/371202}. Retrieval identifiers and hashes accompany the certificate.
\bibitem{patel} Shivam Patel, counterexample reply to MathDB 371202, 20 August 2026. The reply identifier and source retrieval hash are in the provenance record. The cone/count/limit argument is prior work; our primary-source qualification corrects one sampling-attribution assertion in that reply.
\end{thebibliography}
\end{document}
